\section{十二、EC 与 BMC：主机之外仍在运行的控制计算机}

主 CPU 关机、卡死或尚未释放复位时，平台仍要处理电源键、风扇、温度、充电、机箱状态和远程维护。笔记本常由 EC（Embedded Controller，嵌入式控制器）承担键盘、电池与电源时序；服务器常由 BMC（Baseboard Management Controller，基板管理控制器）承担传感器、风扇、事件日志、远程控制台和带外管理网络。它们本身是带处理器、RAM、Flash、定时器和外设总线的小型计算机，并由待机电源供电。

\begin{pdfdiagram}[x=1cm,y=1cm]
  \node[hw state, minimum width=2.3cm] (sensor) at (1.2,4.0) {物理传感器\\温度、电压、电流};
  \node[hw block, minimum width=2.5cm] (sideband) at (4.1,4.0) {带外总线\\I2C/PMBus/PECI};
  \node[hw control, minimum width=2.5cm] (bmc) at (7.1,4.0) {EC/BMC\\待机电源下运行};
  \node[hw memory, minimum width=2.5cm] (log) at (10.1,4.0) {事件与遥测\\时间、对象、阈值};
  \node[hw state, minimum width=2.5cm] (remote) at (13.1,4.0) {带外管理\\告警、控制台、升级};
  \draw[hw bus] (sensor) -- (sideband);
  \draw[hw bus] (sideband) -- (bmc);
  \draw[hw bus] (bmc) -- (log);
  \draw[hw bus] (log) -- (remote);

  \node[hw state, minimum width=2.3cm] (watch) at (1.2,1.0) {策略输入\\watchdog、按钮、命令};
  \node[hw control, minimum width=2.5cm] (policy) at (4.1,1.0) {安全策略\\权限、迟滞、超时};
  \node[hw control, minimum width=2.5cm] (seq) at (7.1,1.0) {执行动作\\风扇、限功率、复位};
  \node[hw block, minimum width=2.5cm] (host) at (10.1,1.0) {主机平台\\SoC、DRAM、设备};
  \node[hw state, minimum width=2.5cm] (verify) at (13.1,1.0) {恢复验收\\服务与数据不变量};
  \draw[hw bus] (watch) -- (policy);
  \draw[hw bus] (policy) -- (seq);
  \draw[hw bus] (seq) -- (host);
  \draw[hw bus] (host) -- (verify);
  \draw[hw bus] (bmc.south) -- (seq.north);
\end{pdfdiagram}
\diagramnote{带外管理的观测链与执行链。上排把物理传感器经带外总线送入 EC/BMC，并形成不依赖主机操作系统的事件记录；下排把 watchdog、按钮或远程命令经过权限和超时策略转换成风扇、限功率或复位动作，最后仍要验证主机服务与数据状态。}

\topic{待机电源让平台在“关机”时仍保留一条控制路径}

ATX 平台的待机电源、移动设备的 always-on rail 或服务器管理电源域可以在主处理器电源关闭时继续给 EC/BMC、实时时钟和少量唤醒逻辑供电。电源键首先成为管理控制器的输入，而不是直接切断主电源。控制器检查电池、适配器、温度、机盖或机箱状态，再按依赖顺序启用稳压器、参考时钟和复位域。

关机同样不是立刻断电。操作系统完成文件系统和设备关闭后，通过固件接口请求目标电源状态；EC/BMC 或电源管理器等待允许条件，关闭主电源轨，并保留能响应唤醒源的最小域。若主机无响应，watchdog 可请求受控复位或强制断电，但强制动作可能丢失未持久化数据，策略必须区分“服务无响应”和“立即危及器件安全”。

\topic{带外总线提供独立观测，但仍有共享故障点}

BMC 可通过 I2C/SMBus、PMBus、PECI 或厂商侧带接口读取温度传感器、VRM 遥测、风扇转速、FRU 标识和处理器状态。它不依赖主机 PCIe、主存或操作系统驱动，因此主机崩溃后仍可能保存证据。但 BMC 与主机可能共享传感器、Flash、电源、时钟或总线复用器；共享器件故障会同时破坏被监控对象和观测路径。

传感器值必须带对象身份、采样时间、有效状态和单位。总线超时不能直接当成“温度为 0”，校准缺失不能与真实低温混为一谈。控制策略要为 missing、stale、out-of-range 和 bus error 定义不同动作；例如关键温度传感器失联时进入保守限功率或关机，而不是沿用最后一次低温读数无限运行。

\topic{EC/BMC 固件也是安全启动和更新链的一部分}

带外控制器能够复位主机、访问平台 Flash、读取遥测并连接管理网络，权限通常高于普通主机软件。它需要独立的固件签名验证、防回滚、调试锁定、密钥保护和 A/B 更新。若只保护主 CPU 启动链而允许 BMC 加载任意固件，攻击者仍可通过电源时序、Flash 写入或管理网络控制整机。

管理网络应与业务网络和不可信设备隔离，远程命令要经过身份认证、授权和审计。恢复出厂设置不能恢复已经撤销的旧密钥或解锁量产调试口。BMC 向主机提供的虚拟介质、远程控制台或主机接口也要限制 DMA 与缓冲范围，避免“带外”成为绕过 IOMMU 和操作系统权限的另一条数据通道。

\topic{主机复位后，管理控制器要区分旧状态与新一代启动}

BMC 记录一次复位原因时，应锁存触发源、主机电源状态、watchdog 阶段和关键传感器快照。主机重新启动后，固件读取这些记录并确认已经消费，再由明确协议清除。若新启动一开始就清空所有状态，最有价值的故障证据会在诊断前丢失；若永不清除，旧告警又会被误归因到新故障。

管理控制器自身也可能卡死。独立 watchdog、只读恢复镜像或更小的电源时序状态机可以重启 BMC，但要避免在 BMC 复位时无条件切断正常主机。恢复后先重建传感器总线、风扇安全策略和事件代际，再恢复远程管理服务；主机电源动作必须根据当前真实状态重新判定，不能重放复位前队列中的旧命令。

\begin{longtable}{L{0.17\linewidth}L{0.28\linewidth}L{0.31\linewidth}L{0.14\linewidth}}
\toprule
管理对象 & 必须保存的状态 & 安全默认动作 & 关键故障 \\
\midrule
温度/电源传感器 & 对象、单位、时间、有效性、阈值 & 失联时限功率或进入安全态 & 旧值被当成当前值 \\\tablerowrule
风扇与泵 & 目标转速、实际转速、控制代际 & 提高冗余风扇并告警 & 共用供电导致同时停转 \\\tablerowrule
watchdog & 阶段、最后进展点、超时次数 & 先保存证据再选择复位域 & 复位循环清除根因 \\\tablerowrule
平台 Flash & 当前槽、候选槽、回滚计数 & 保持已验证旧槽可启动 & BMC 与主机同时改写 \\\tablerowrule
远程管理 & 身份、权限、命令、审计时间 & 拒绝未认证的高权限动作 & 管理面绕过主机隔离 \\
\bottomrule
\end{longtable}

\section{十三、从 FIT 和 MTTR 算到平台可用性}

“器件可靠”与“服务可用”不是同一个指标。可靠性讨论一段时间内不发生失效的概率；可用性还把检测、切换、维修和恢复时间计入。一个部件很少失效，但每次维修数天，可能比一个偶尔失效却能在数秒内自动切换的系统造成更多停机。平台设计需要同时记录失效率、故障覆盖率、平均检测时间、平均修复时间和恢复后的正确性。

\topic{FIT 与 MTBF 给出不同方向的同一数量级}

FIT（Failures In Time）表示每 $10^9$ 器件小时的预期失效数。在近似恒定失效率模型下，若部件为 100 FIT，则失效率约为 $100/10^9=10^{-7}$ 每小时，MTBF（Mean Time Between Failures）约为 $10^7$ 小时。它是大量同类器件的统计量，不表示某一器件会在运行到第 $10^7$ 小时时准时失效，也不适合直接描述早期失效和磨损期明显变化的浴盆曲线。

串联系统只要任一关键部件失效就停止服务，在独立、低失效率近似下，总失效率约为各部件失效率之和。十个各 100 FIT 的关键部件串联，平台关键路径约为 1000 FIT，近似 MTBF 为 $10^6$ 小时。这个结果说明增加部件即使每个都很可靠，也会增加必须被检测和隔离的故障来源。

\topic{MTTR 决定失效发生后停多久}

稳态可用性可近似写为
\[
A=\frac{\mathrm{MTBF}}{\mathrm{MTBF}+\mathrm{MTTR}}.
\]
若某服务 MTBF 为 100000 小时，MTTR 为 4 小时，则 $A\approx99.996\%$，折算每年期望停机约 $8760(1-A)\approx0.35$ 小时，即约 21 分钟。把 MTTR 从 4 小时降到 10 分钟，即使器件失效率不变，也能显著减少停机。

MTTR 不只是更换硬件所需时间。它包含检测、定位、切换、修复、重新同步和验收。错误已被 watchdog 检出但等待人工确认一小时，这一小时属于不可用时间；设备复位后立即出现，却还要重建 RAID、Cache 或连接状态，也尚未恢复完整服务。

\topic{冗余收益取决于独立性和故障覆盖率}

两个并联副本若真正独立，且任一副本都能承担服务，同时失效概率可以远低于单副本。但系统还必须检测主副本已失效、正确隔离错误输出并完成切换。若故障检测覆盖率不足，坏副本可能继续提供错误数据；若两副本共用电源、时钟、冷却、交换机或同一固件缺陷，共同原因会使独立概率乘法失效。

冗余控制器本身也可能成为单点。主备仲裁需要明确谁有权写共享状态，网络分区时要避免双方同时成为主节点；双写数据要定义一致性和持久化边界；恢复旧副本前要先同步状态，不能把过期数据重新投票为正确。可靠性框图应把电源、管理、软件和运维动作画进故障域，而不是只数业务数据副本。

\topic{降级模式要进入容量和时延预算}

冗余部件失去一个后，系统可能仍在线，但剩余链路、风扇、内存通道或存储副本要承担更高负载。正常状态只使用 50\% 容量，并不保证降级后仍满足尾延迟；故障切换期间的缓存冷启动、路由收敛和数据重建还会进一步占用资源。可用性验收应在故障注入后施加目标负载，验证数据正确、容量足够、超时假设仍成立。

\begin{longtable}{L{0.18\linewidth}L{0.24\linewidth}L{0.30\linewidth}L{0.18\linewidth}}
\toprule
指标 & 回答的问题 & 容易遗漏的条件 & 适合的证据 \\
\midrule
FIT/失效率 & 单位时间预期出现多少故障 & 生命周期阶段与环境应力 & 器件统计、温度和负载 \\\tablerowrule
MTBF & 两次可修复故障平均间隔 & 不是单器件寿命倒计时 & 现场故障间隔 \\\tablerowrule
MTTR & 从失效到完整恢复需要多久 & 检测、同步和验收也计入 & 恢复时间线 \\\tablerowrule
可用性 & 长期有多少时间能提供服务 & 服务等级与降级容量定义 & SLO 与停机窗口 \\\tablerowrule
覆盖率 & 故障能否被检测并正确隔离 & 潜伏错误与误纠正 & 故障注入结果 \\\tablerowrule
共同原因 & 哪些副本会被同一事件同时破坏 & 电源、冷却、固件和运维共享 & 故障域与依赖图 \\
\bottomrule
\end{longtable}
